\documentclass[11pt]{article}
\usepackage[T1]{fontenc}
\usepackage{lmodern,amsmath,amssymb,amsthm,mathtools}
\usepackage[a4paper,margin=28mm]{geometry}
\usepackage{xurl}
\usepackage[hidelinks,breaklinks]{hyperref}
\usepackage{microtype}
\DeclareMathOperator{\per}{per}
\newtheorem{theorem}{Theorem}
\newtheorem{lemma}[theorem]{Lemma}
\newtheorem{corollary}[theorem]{Corollary}
\title{The sharp four-row permanent--determinant inequality and its complete equality classification}
\author{Yongxian Zhang\\[2pt]
\small School of Computer Science and Engineering\\
\small South China University of Technology\\
\small\texttt{mxymmxym1@gmail.com}\\
\small ORCID: 0009-0000-3864-3536}
\date{13 October 2026}
\begin{document}
\maketitle
\begin{abstract}
For every complex $4\times4$ matrix $A$ and every real $c\ge0$, we prove
$|\operatorname{per} A|+c|\det A|\le\max\{3/2,1+c\}\prod_i\|A_{i,*}\|_2$.
The coefficient is optimal. We classify equality, including zero rows and
the transition $c=1/2$: the nondegenerate extremizers are the rank-one
matrices with an equimodular common row direction, or the monomial matrices,
with the relevant class selected by $c$. The exact real-coefficient norm
of $\operatorname{per} A+t\det A$ is also determined. These principal results have
complete Lean proofs for the actual matrix permanent and determinant.
The written proof uses two-row identities and Laplace expansion.
\end{abstract}

\section{Definitions and principal results}
For $A\in\mathbb C^{4\times4}$, let
\[
 \per A=\sum_{\sigma\in S_4}\prod_{i=1}^4 A_{i,\sigma(i)},\qquad
 \det A=\sum_{\sigma\in S_4}\operatorname{sgn}(\sigma)
                         \prod_{i=1}^4 A_{i,\sigma(i)},
 \qquad R(A)=\prod_{i=1}^4\Big(\sum_{j=1}^4|A_{ij}|^2\Big)^{1/2}.
\]
A matrix is \emph{monomial} if there is a permutation $\sigma$ such that
$A_{i,\sigma(i)}\ne0$ for every $i$ and all its other entries vanish.
A nonzero \emph{flat rank-one} matrix has the form
$A_{ij}=u_i v_j$, where every $u_i$ and $v_j$ is nonzero and
$|v_1|=\cdots=|v_4|$.

\begin{theorem}\label{thm:main}
For every $A\in\mathbb C^{4\times4}$ and every real $c\ge0$,
\begin{equation}\label{eq:main}
 |\per A|+c|\det A|\le M(c)R(A),\qquad
 M(c)=\max\{3/2,1+c\}.
\end{equation}
Equality holds if and only if at least one of the following applies:
\begin{enumerate}
\item $A$ has a zero row;
\item $A$ is flat rank one and $c\le1/2$;
\item $A$ is monomial and $c\ge1/2$.
\end{enumerate}
The constant $M(c)$ is attained and cannot be decreased.
\end{theorem}

The permanent endpoint agrees with the sharp Hadamard-type permanent
inequality of Carlen, Lieb and Loss \cite{CLL}. The joint bound and its
transition are proved below directly; their proof does not assume the
sharp permanent bound. We make no exhaustive historical-priority claim.

\begin{corollary}\label{cor:pencil}
For every real $t$,
\[
 \sup_{A:\,R(A)>0}\frac{|\per A+t\det A|}{R(A)}
       =\max\{3/2,1+|t|\}.
\]
\end{corollary}
\begin{proof}
The triangle inequality and Theorem~\ref{thm:main}, with $c=|t|$, give
the upper bound. The matrix whose entries are all $1/2$ has $R=1$,
$\per=3/2$ and $\det=0$. An even permutation matrix if $t\ge0$, and
an odd permutation matrix if $t<0$, has $R=1$ and objective $1+|t|$.
Thus the larger endpoint is attained.
\end{proof}

\section{The two-row calculation}
For unit vectors $a,b\in\mathbb C^4$, define
\[
 S(a,b)=\sum_{j<k}|a_jb_k+a_kb_j|^2,\qquad
 W(a,b)=\sum_{j<k}|a_jb_k-a_kb_j|^2.
\]
Write $\eta=|\sum_j a_j\overline{b_j}|^2$,
$\tau=\sum_j|a_j|^2|b_j|^2$ and $\Delta=\tau-\eta/4$.
Cauchy--Schwarz gives $0\le\eta\le1$ and $\Delta\ge0$.
Expansion of the squares gives the exact identities
\begin{equation}\label{eq:pair}
 S=1+\eta-2\tau,\qquad W=1-\eta,\qquad
 F_c:=S+cW=1+c+(1/2-c)\eta-2\Delta.
\end{equation}
In the second use of Cauchy--Schwarz, $\eta\le4\tau$ is applied
to the four numbers $a_j\overline{b_j}$. Its equality condition is
\begin{equation}\label{eq:product}
 \Delta=0\quad\Longleftrightarrow\quad
 a_j\overline{b_j}\text{ is independent of }j.
\end{equation}
Consequently
\begin{equation}\label{eq:pairbound}
 0\le F_c(a,b)\le M(c).
\end{equation}
If $c<1/2$, equality in the upper bound requires $\eta=1$ and
$\Delta=0$. If $c>1/2$, it requires $\eta=\Delta=0$.
At $c=1/2$, it requires exactly $\Delta=0$.

\section{The four-row inequality and sharpness}
If a row is zero, both multilinear expressions vanish, so
\eqref{eq:main} is an equality. Otherwise divide each row by its positive
Euclidean norm. It suffices to prove the assertion for unit rows
$a,b,d,e$. For a two-column subset $J=\{j,k\}$, $j<k$, put
\[
 p_{ab}(J)=a_jb_k+a_kb_j,\qquad
 w_{ab}(J)=a_jb_k-a_kb_j.
\]
The permanent and determinant Laplace expansions along the first two rows
are
\[
 \per A=\sum_{|J|=2}p_{ab}(J)p_{de}(J^c),\qquad
 \det A=\sum_{|J|=2}\varepsilon_J w_{ab}(J)w_{de}(J^c),
\]
where $\varepsilon_J\in\{-1,1\}$ is the column-shuffle sign.
The complement map permutes the six two-column subsets. Hence
Cauchy--Schwarz, first on these six terms and then on two real coordinates,
gives
\begin{align}
 |\per A|+c|\det A|
 &\le\sqrt{S(a,b)S(d,e)}+c\sqrt{W(a,b)W(d,e)}\notag\\
 &\le\sqrt{F_c(a,b)F_c(d,e)}\le M(c).\label{eq:laplace}
\end{align}
The second step also holds at $c=0$; for $c>0$ it is the ordinary
Cauchy--Schwarz inequality applied to
$(\sqrt S,\sqrt{cW})$ for the two row pairs. Restoring the row norms
proves~\eqref{eq:main}.

The matrix with all entries $1/2$ has unit rows and objective $3/2$.
Every permutation matrix has unit rows and objective $1+c$. These
examples prove optimality and attainment of $M(c)$ for every $c\ge0$.

\section{Complete equality classification}
Assume all rows are nonzero and normalize them as above. Permuting rows
preserves the absolute values of the permanent and determinant. Thus
\eqref{eq:laplace} applies to every division of the four row indices into
two pairs. If its left side equals $M(c)$, both pair factors must equal
$M(c)$: each is nonnegative and at most the positive number $M(c)$.
Every unordered pair occurs in such a division, so all six pairs satisfy
the equality conditions following~\eqref{eq:pairbound}.

If $c<1/2$, every pair has $\eta=1$, so all unit rows are proportional.
For a pair of proportional unit rows, $\Delta=0$ says
$\sum_j|a_j|^4=1/4$. Since $\sum_j|a_j|^2=1$, equality in
Cauchy--Schwarz forces $|a_j|^2=1/4$ for all four columns.
The matrix is therefore flat rank one.

If $c>1/2$, every pair has $\eta=\Delta=0$, hence $\tau=0$.
The row supports are pairwise disjoint. Four nonempty pairwise disjoint
subsets of four columns must each be a distinct singleton, which gives
exactly a monomial matrix.

At $c=1/2$, each pair satisfies~\eqref{eq:product}. Denote its constant
by $z_{ih}$, so
\begin{equation}\label{eq:z}
 a_{ij}\overline{a_{hj}}=z_{ih}\quad(i\ne h)
\end{equation}
for every column $j$. If the four row supports are pairwise disjoint, the
same argument gives a monomial matrix. Otherwise two rows meet in one
column, so their $z_{ih}$ is nonzero and~\eqref{eq:z} forces both rows
to have full support. Each remaining nonzero row meets one of these full
rows and therefore also has full support. All the $z_{ih}$ are now
nonzero. Choose distinct rows $i,h,k$. For every column $j$,
\[
 |a_{ij}|^2=\frac{|z_{ih}|\,|z_{ik}|}{|z_{hk}|}.
\]
This is independent of $j$. Unit normalization makes every entry of
every row have modulus $1/2$. Relative to a fixed reference row, division
in~\eqref{eq:z} then makes each row a scalar multiple of that row.
The matrix is flat rank one.

Conversely, for a flat rank-one matrix $A_{ij}=u_i v_j$ with common
modulus $|v_j|=r>0$,
\[
 |\per A|=24\Big(\prod_i|u_i|\Big)r^4
          =\tfrac32 R(A),\qquad\det A=0.
\]
A monomial matrix satisfies $|\per A|=|\det A|=R(A)$.
These formulas give equality precisely at the stated weights.
Together with the zero-row case this finishes the proof of
Theorem~\ref{thm:main}.

\section{Formal verification and scope}
The repository \url{https://github.com/mxym/math} contains the fixed
Lean projects \nolinkurl{formalizations/four-row-permanent-tradeoff} and
\nolinkurl{formalizations/four-row-complete-equality}. At repository
revision \nolinkurl{07559d229} the latter exports
\nolinkurl{FourRowTradeoff.sharp_four_row_complete}; its matrix classes
are \texttt{IsFlatRankOne} and \texttt{IsMonomial}. It covers the actual
complex matrix inequality, both directions of equality, zero rows and
the critical weight. The former project supplies optimality, attainment
and the real-coefficient pencil norm. The complete-equality project's
\texttt{FORMALIZATION\_MAP.md} records the exact correspondence.

Lean is pinned to 4.34.1 (commit
\nolinkurl{5045d0056413266e57c625dcd7c365b10e377c52}); Mathlib is pinned to
\nolinkurl{d13f23b723b8a846827a245b89c10fc7d3f11612}.
The archived verification completed on 9 October 2026 and reports a
fresh build of all 12 owned modules and a replay of 112 named roots with
a 15,593-declaration dependency closure into an empty kernel at trust
level zero. The only logical axioms are \texttt{propext},
\texttt{Classical.choice} and \texttt{Quot.sound}; two invalid-proof
controls were rejected. Reproduction commands and source hashes are in
that project's README and \texttt{evidence/v2/VERIFICATION.json}.
These are archived results; preparation of this reading edition does
not claim a new Lean execution.

This paper consolidates the principal fully formalized results of the
repository note \texttt{notes/four-row-permanent-tradeoff/PAPER.md}.
Its additional rectangular, quantitative, convex-power and tensor
extensions are outside this paper's formalization claim. No external
human peer-review acceptance or mathematical-priority certification is
claimed. Research and manuscript preparation used OpenAI AI assistants;
the author is responsible for the stated claims. There was no external
funding. Original material is all rights reserved, subject to existing
third-party licenses and notices.

\begin{thebibliography}{9}
\bibitem{CLL} E. Carlen, E. H. Lieb and M. Loss,
\emph{An inequality of Hadamard type for permanents},
Methods and Applications of Analysis \textbf{13} (2006), 1--18.
\href{https://doi.org/10.4310/MAA.2006.v13.n1.a1}{doi:10.4310/MAA.2006.v13.n1.a1}.
\bibitem{repo} Y. Zhang,
\emph{Sharp four-row permanent--determinant tradeoff and parity tensor norms},
research note and Lean source projects, 2026.
\url{https://github.com/mxym/math/tree/07559d229/notes/four-row-permanent-tradeoff}.
\end{thebibliography}
\end{document}
